\section{Grundlagen der Flüssigkeitskühlung}
\label{sec:grundlagen_kuehlung}

\subsection{Wärmeentstehung und Kühlkonzept}
\label{subsec:waermeentstehung}

Im Betrieb eines elektrisch angetriebenen Formula-Student-Fahrzeugs entstehen in
mehreren Komponenten thermische Verluste. Maßgeblich sind die Elektromotoren und
die Wechselrichter, die einen Teil der aufgenommenen elektrischen Leistung als
Wärme an ihre Umgebung abgeben. Wird diese Wärme nicht ausreichend abgeführt,
steigen die Bauteiltemperaturen an. Die Folgen sind eine Leistungsreduktion
(Derating), im ungünstigen Fall die Abschaltung oder die Beschädigung der
Komponenten.

Bei der Flüssigkeitskühlung nimmt ein Kühlmittel die Verlustwärme direkt an der
Komponente auf und transportiert sie zu einem Wärmetauscher (Radiator). Dort
wird die Wärme an die Umgebungsluft abgegeben, und das abgekühlte Kühlmittel
strömt zurück zu den Komponenten. Es entsteht ein geschlossener Kühlkreislauf
(Abbildung~\ref{fig:kuehlkreislauf_prinzip}).

% --- Bildvorschlag 1 ---
\begin{figure}[htb]
  \centering
  % \includegraphics[width=0.8\textwidth]{bilder/kuehlkreislauf_prinzip}
  \fbox{\parbox{0.8\textwidth}{\centering\vspace{3cm}Platzhalter: Prinzipskizze
  Kühlkreislauf (Pumpe, Motoren, Inverter, Radiator, Lüfter)\vspace{3cm}}}
  \caption{Prinzipieller Aufbau eines geschlossenen Flüssigkeitskühlkreislaufs}
  \label{fig:kuehlkreislauf_prinzip}
\end{figure}

\subsection{Wärmetransport im Kühlmittel}
\label{subsec:waermetransport}

Die vom Kühlmittel abgeführte Wärmeleistung ergibt sich aus dem Massenstrom
und der Temperaturänderung zwischen Ein- und Austritt der Komponente:
\begin{equation}
  \dot{Q} = \dot{m} \cdot c_p \cdot \Delta T
  \label{eq:waermeleistung}
\end{equation}
mit
\begin{itemize}
  \item $\dot{Q}$ -- abgeführte Wärmeleistung in \si{\watt},
  \item $\dot{m}$ -- Massenstrom des Kühlmittels in \si{\kilo\gram\per\second},
  \item $c_p$ -- spezifische Wärmekapazität des Kühlmittels in
        \si{\joule\per\kilo\gram\per\kelvin},
  \item $\Delta T$ -- Temperaturdifferenz zwischen Ein- und Austritt in \si{\kelvin}.
\end{itemize}

Gemessen wird in der Praxis der Volumenstrom $\dot{V}$. Der Massenstrom folgt
über die Dichte $\rho$ des Kühlmittels:
\begin{equation}
  \dot{m} = \rho \cdot \dot{V}
  \label{eq:massenstrom}
\end{equation}
Gleichung~\eqref{eq:waermeleistung} zeigt, dass sich eine gegebene Wärmeleistung
über einen großen Volumenstrom bei kleinem $\Delta T$ oder über einen kleinen
Volumenstrom bei großem $\Delta T$ abführen lässt. Ein kleines $\Delta T$ ist
günstig, weil die Komponenten dann gleichmäßiger und auf niedrigerem
Temperaturniveau gekühlt werden. Der Volumenstrom ist jedoch durch die Pumpe und
den Druckverlust des Kreislaufs begrenzt (vgl.
Abschnitt~\ref{subsec:hydraulik}). Das gilt besonders bei hintereinander
geschalteten Komponenten, denn dort erwärmt sich das Kühlmittel entlang des
Strömungswegs, und die letzte Komponente wird mit dem wärmsten Kühlmittel
versorgt.

Gegenüber einer direkten Luftkühlung hat die Flüssigkeitskühlung einen
wesentlichen Vorteil. Bezogen auf das Volumen speichert Wasser deutlich mehr
Wärme als Luft: Die volumetrische Wärmekapazität $\rho \cdot c_p$ liegt bei
Wasser etwa um den Faktor \num{3500} über der von Luft. Dadurch lassen sich auch
bei kompakter Bauweise hohe Wärmeströme abführen. Die Kennwerte des verwendeten
Kühlmittels sind in Tabelle~\ref{tab:stoffwerte} angegeben.

\begin{table}[htb]
  \centering
  \caption{Stoffwerte von destilliertem Wasser bei Umgebungsdruck}
  \label{tab:stoffwerte}
  \begin{tabular}{|l|c|c|c|}
    \hline
    Größe & \multicolumn{2}{c|}{Wert} & Einheit \\
    \cline{2-3}
          & $20\,^{\circ}\mathrm{C}$ & $60\,^{\circ}\mathrm{C}$ & \\
    \hline
    Dichte $\rho$              & 998  & 983  & $\mathrm{kg/m^3}$ \\
    \hline
    Spez. Wärmekapazität $c_p$ & 4182 & 4185 & $\mathrm{J/(kg\,K)}$ \\
    \hline
    Dyn. Viskosität $\eta$     & $1{,}00 \cdot 10^{-3}$ & $0{,}47 \cdot 10^{-3}$ & $\mathrm{Pa\,s}$ \\
    \hline
  \end{tabular}
\end{table}

\subsection{Hydraulische Auslegung}
\label{subsec:hydraulik}

Neben dem Wärmetransport bestimmen die hydraulischen Eigenschaften die Auslegung
des Kühlsystems. In Schläuchen, Kupplungen, Radiator, Motoren und Inverter
entstehen Druckverluste, die die Pumpe überwinden muss. Der Gesamtdruckverlust
einer Reihenschaltung ist die Summe der Einzelverluste:
\begin{equation}
  \Delta p_\mathrm{ges} = \sum_i \Delta p_i
  \label{eq:dp_ges}
\end{equation}
Bei turbulenter Strömung wächst der Druckverlust einer Komponente näherungsweise
quadratisch mit dem Volumenstrom. Er lässt sich über einen
Widerstandskoeffizienten $R_i$ beschreiben:
\begin{equation}
  \Delta p_i = R_i \cdot \dot{V}^{2}
  \label{eq:dp_komponente}
\end{equation}
Dieser Zusammenhang ist die Anlagenkennlinie. Zusammen mit der Pumpenkennlinie
ergibt sich der Betriebspunkt des Systems (Abbildung~\ref{fig:pumpe_anlage}):
Die Pumpe stellt genau den Volumenstrom ein, bei dem der von ihr gelieferte
Druck dem Druckverlust der Anlage entspricht. Bei Reihenschaltung addieren sich
die Widerstände. Bei Parallelschaltung teilt sich der Volumenstrom auf die
Zweige auf, und der Druckverlust ist in allen Zweigen gleich. Eine Änderung am
Layout, etwa zusätzliche Komponenten oder eine andere Schlauchführung,
verschiebt daher den Betriebspunkt und damit den abführbaren Wärmestrom nach
Gleichung~\eqref{eq:waermeleistung}.

% --- Bildvorschlag 2 ---
\begin{figure}[htb]
  \centering
  % \includegraphics[width=0.7\textwidth]{bilder/pumpe_anlagenkennlinie}
  \fbox{\parbox{0.7\textwidth}{\centering\vspace{3cm}Platzhalter: Pumpenkennlinie
  und Anlagenkennlinie $\Delta p(\dot{V})$ mit Betriebspunkt\vspace{3cm}}}
  \caption{Pumpen- und Anlagenkennlinie mit resultierendem Betriebspunkt}
  \label{fig:pumpe_anlage}
\end{figure}

Für die Auslegung muss die Pumpe somit den benötigten Volumenstrom gerade bei
dem im System auftretenden Gesamtdruckverlust liefern. Die Widerstandswerte der
Einzelkomponenten des Fahrzeugs wurden in dieser Arbeit messtechnisch bestimmt
(Abschnitt~\ref{sec:hydraulik_vermessung}) und bilden die Grundlage für die
Simulation des Kreislaufs.

\section{Thermische Anforderungen}
\label{sec:thermische_anforderungen}

Das Kühlsystem muss gewährleisten, dass alle thermisch kritischen Komponenten
während des gesamten Fahrzeugbetriebs innerhalb ihrer zulässigen
Temperaturbereiche betrieben werden. Zu betrachten sind insbesondere die
Elektromotoren und die Wechselrichter.

\paragraph{Kreislaufaufteilung.}
Das Fahrzeug besitzt zwei getrennte Kühlkreisläufe
(Abbildung~\ref{fig:kreislaufaufteilung}):
\begin{itemize}
  \item \textbf{Kreislauf 1:} Versorgung der Motoren der Hinterachse.
        Die maximal zulässige Kühlmitteltemperatur beträgt \SI{90}{\celsius}.
  \item \textbf{Kreislauf 2:} Das Kühlmittel durchströmt zunächst die
        Wechselrichter und wird anschließend auf die Motoren der Vorderachse
        verteilt. Wegen der geringeren zulässigen Betriebstemperatur der
        Wechselrichter wird für diesen Kreislauf eine maximale
        Kühlmitteltemperatur von \SI{60}{\celsius} angesetzt.
\end{itemize}

% --- Bildvorschlag 1 ---
\begin{figure}[htb]
  \centering
  % \includegraphics[width=0.85\textwidth]{bilder/kreislaufaufteilung}
  \fbox{\parbox{0.85\textwidth}{\centering\vspace{3cm}Platzhalter: Schaltbild
  beider Kreisläufe (Kreislauf 1: Pumpe, Radiator, Motoren HA; Kreislauf 2:
  Pumpe, Radiator, Inverter, Motoren VA) mit Temperaturgrenzen\vspace{3cm}}}
  \caption{Aufteilung des Kühlsystems in zwei getrennte Kreisläufe}
  \label{fig:kreislaufaufteilung}
\end{figure}

\paragraph{Auslegungsgrößen.}
Neben den zulässigen Temperaturen ist die abzuführende Verlustleistung die
wesentliche Auslegungsgröße. Das Kühlsystem muss auch unter ungünstigen
Bedingungen eine ausreichende Wärmeabfuhr sicherstellen. Als kritischer
Auslegungsfall wird eine hohe Umgebungstemperatur von bis zu
\SI{40}{\celsius} angesetzt. Die erforderliche Kühlleistung folgt aus der
thermischen Belastung der Komponenten (Tabelle~\ref{tab:thermische_anforderungen}).

\begin{table}[htb]
  \centering
  \caption{Thermische Auslegungsgrößen der beiden Kühlkreisläufe}
  \label{tab:thermische_anforderungen}
  \begin{tabular}{|l|c|c|}
    \hline
    Größe & Kreislauf 1 & Kreislauf 2 \\
    \hline
    Gekühlte Komponenten & Motoren HA & Inverter, Motoren VA \\
    \hline
    Max. Kühlmitteltemperatur $T_\mathrm{K,max}$ & $90\,^{\circ}\mathrm{C}$ & $60\,^{\circ}\mathrm{C}$ \\
    \hline
    Max. Umgebungstemperatur $T_\mathrm{U,max}$ & $40\,^{\circ}\mathrm{C}$ & $40\,^{\circ}\mathrm{C}$ \\
    \hline
    Treibende Temperaturdifferenz $T_\mathrm{K,max}-T_\mathrm{U,max}$ & $50\,\mathrm{K}$ & $20\,\mathrm{K}$ \\
    \hline
    Abzuführende Verlustleistung $\dot{Q}_\mathrm{V}$ & \textbf{??}\,W & \textbf{??}\,W \\
    \hline
  \end{tabular}
\end{table}

Die Wärmeabgabe des Radiators an die Umgebungsluft ist näherungsweise
proportional zur Differenz zwischen Kühlmittel-Eintrittstemperatur und
Lufteintrittstemperatur:
\begin{equation}
  \dot{Q}_\mathrm{Rad} \approx \varepsilon \cdot C_\mathrm{min} \cdot
  \left(T_\mathrm{K,ein} - T_\mathrm{L,ein}\right)
  \label{eq:radiator}
\end{equation}
mit der Effektivität $\varepsilon$ des Wärmetauschers und dem kleineren der
beiden Wärmekapazitätsströme $C_\mathrm{min}$ (in der Regel der Luftseite). Bei
hohen Außentemperaturen sinkt die treibende Temperaturdifferenz. Bei
Kreislauf~2 beträgt sie im Auslegungsfall nur \SI{20}{\kelvin} gegenüber
\SI{50}{\kelvin} bei Kreislauf~1. Um dieselbe Wärmeleistung abzuführen, ist daher
ein entsprechend größerer Luftmassenstrom bzw. eine größere Kühlerfläche
erforderlich. Der Wärmestrom auf der Luftseite ergibt sich analog zu
Gleichung~\eqref{eq:waermeleistung}:
\begin{equation}
  \dot{Q}_\mathrm{Rad} = \dot{m}_\mathrm{L} \cdot c_{p,\mathrm{L}} \cdot
  \Delta T_\mathrm{L}
  \label{eq:luftseite}
\end{equation}

Ziel der thermischen Auslegung ist es, Kühlmittel- und Luftmassenstrom so zu
bestimmen, dass die zulässigen Bauteiltemperaturen auch unter den definierten
Worst-Case-Bedingungen nicht überschritten werden.

% --- Bildvorschlag 2 ---
\begin{figure}[htb]
  \centering
  % \includegraphics[width=0.7\textwidth]{bilder/temperaturgrenzen}
  \fbox{\parbox{0.7\textwidth}{\centering\vspace{2.5cm}Platzhalter: Diagramm der
  Temperaturniveaus (Umgebung \SI{40}{\celsius}, Grenze Kreislauf 2
  \SI{60}{\celsius}, Grenze Kreislauf 1 \SI{90}{\celsius}) mit
  Temperaturdifferenzen\vspace{2.5cm}}}
  \caption{Zulässige Kühlmitteltemperaturen und verbleibende Temperaturdifferenz zur Umgebung}
  \label{fig:temperaturgrenzen}
\end{figure}


\section{Randbedingungen des Formula-Student-Fahrzeugs}
\label{sec:randbedingungen}

Neben den thermischen und konstruktiven Anforderungen sind bei der Entwicklung
des Kühlsystems die Vorgaben des Formula-Student-Regelwerks einzuhalten. Das
Regelwerk der Formula Student Germany definiert Anforderungen an das
Kühlmedium, die Temperaturbeständigkeit, die Dichtheit und die konstruktive
Ausführung des Kühlsystems~\cite{fsg_rules}. Tabelle~\ref{tab:regeln} fasst die
für diese Arbeit relevanten Regeln zusammen.

\begin{table}[htb]
  \centering
  \caption{Relevante Regelwerksanforderungen an das Kühlsystem}
  \label{tab:regeln}
  \begin{tabular}{|l|p{0.78\textwidth}|}
    \hline
    Regel & Anforderung \\
    \hline
    T7.2.2   & Als Kühlmedium sind ausschließlich Wasser, Luft oder Öl zulässig. \\
    \hline
    T7.2.3   & Wassergekühlte Systeme benötigen eine hitzebeständige, starr
               befestigte Abdeckung (dauerhaft mindestens $100\,^{\circ}\mathrm{C}$);
               Ausnahme: außenliegende Radnabenmotoren und deren Kühlleitungen. \\
    \hline
    T7.2.4   & Das Austreten von Flüssigkeit muss verhindert werden. \\
    \hline
    T7.2.6   & Bei vorhandenen Entlüftungen ist ein Auffangbehälter mit mindestens
               $10\,\%$ des Fluidvolumens oder $100\,\mathrm{ml}$
               (der größere Wert) erforderlich. \\
    \hline
    T7.2.7   & Alle Komponenten des Motor-Kühlsystems inkl. Befestigungen sind für
               mindestens $100\,^{\circ}\mathrm{C}$ bzw. die maximal zu erwartende
               Mediumstemperatur auszulegen (der höhere Wert gilt). \\
    \hline
    T7.3.5   & Rotierende Komponenten im Stillstand (z.\,B. Kühlerventilatoren)
               benötigen einen Fingerschutz; Perforationen dürfen das Durchdringen
               von Objekten ab $12\,\mathrm{mm}$ Durchmesser nicht zulassen. \\
    \hline
    T11.11.1 & Die Summe der Nennleistungen aller aktiven Luftfördereinrichtungen
               (inkl. Kühlerlüfter) darf $500\,\mathrm{W}$ nicht überschreiten. \\
    \hline
  \end{tabular}
\end{table}

\paragraph{Kühlmedium.}
Nach T7.2.2 sind nur Wasser, Luft oder Öl als Kühlmedium zulässig. Für das
entwickelte Kühlsystem wird daher destilliertes Wasser verwendet.

\paragraph{Fahrerschutz und Temperaturbeständigkeit.}
Wassergeführte Systeme müssen nach T7.2.3 mit einer starr befestigten,
dauerhaft für mindestens \SI{100}{\celsius} ausgelegten Abdeckung versehen
sein. Ausgenommen sind außenliegende Radnabenmotoren und deren Kühlleitungen.
Gemäß T7.2.7 müssen zudem alle Komponenten des Motor-Kühlsystems einschließlich
ihrer Befestigungen (Kühler, Pumpen, Leitungen, Adapter, Anschlüsse) für
mindestens \SI{100}{\celsius} oder die maximal zu erwartende Mediumstemperatur
ausgelegt sein, je nachdem, welcher Wert höher ist.

\paragraph{Dichtheit und Entlüftung.}
T7.2.4 fordert, dass keine Flüssigkeit austreten darf. Die Dichtheit der
Leitungen, Anschlüsse und Verbindungselemente ist damit eine zentrale
Konstruktionsanforderung. Bei vorhandenen Entlüftungen schreibt T7.2.6 einen
Auffangbehälter vor. Das Kühlsystem des RP26E/D ist als geschlossenes System
ohne Entlüftungsbehälter ausgeführt. Dadurch entfallen zusätzliche Komponenten
und deren Befestigung, und es wird Gewicht eingespart. Begründung, warum
keine Entlüftung im Sinne von T7.2.6 vorliegt (siehe Abschnitt
\ref{sec:entlueftung})

\paragraph{Schutz rotierender Bauteile.}
Für die Kühlerlüfter ist nach T7.3.5 ein mechanischer Fingerschutz erforderlich,
der Fingerkräften standhält. Perforierte Materialien sind zulässig, sofern
Objekte ab \SI{12}{\milli\meter} Durchmesser nicht hindurchgelangen.

\paragraph{Begrenzung der Lüfterleistung.}
Nach T11.11.1 darf die Summe der Nennleistungen aller aktiven Luftfördereinrichtungen
\SI{500}{\watt} nicht überschreiten. Die Kühlerlüfter fallen ausdrücklich
darunter. Der erreichbare Luftmassenstrom nach Gleichung~\eqref{eq:luftseite}
ist damit nach oben begrenzt. Dies steht im direkten Zielkonflikt zu den
erhöhten Anforderungen bei hoher Umgebungstemperatur
(Abschnitt~\ref{sec:thermische_anforderungen}).

\paragraph{Zusammenfassung.}
Neben den thermischen und hydraulischen Anforderungen bilden die
Regelwerksvorgaben weitere Randbedingungen. Dichtheit, Temperaturbeständigkeit,
Schutz bewegter Bauteile und die Begrenzung der Lüfterleistung sind bereits
in der Konstruktion zu berücksichtigen. Die Ausführung als geschlossenes System
ohne Entlüftungsbehälter reduziert die Systemmasse, ohne das Regelwerk zu
verletzen, sofern keine Entlüftung des wasserführenden Systems erforderlich ist.

% --- Bildvorschlag 3 ---
% Foto/CAD: Fingerschutz am Kühlerlüfter bzw. Abdeckung der Kühlleitungen
% (zeigt die Umsetzung von T7.2.3 / T7.3.5 im Fahrzeug)